\documentclass[11pt]{article}

\usepackage[margin=1in]{geometry}
\usepackage{amsmath}
\usepackage{amssymb}

\title{1PN vs 0PN dephasing}
\author{}
\date{}

\begin{document}

\maketitle

Velcani computed the wrong quantity for a coherent search. The fractional PN
correction to $\Psi(f)$ is small, but the quantity of interest is dephasing. For an
equal-mass $M_c=10^{-3}M_\odot$ binary starting at 20 Hz, the raw accumulated
1PN dephasing by 120 Hz is already $1.42\times10^6$ rad, so a local fractional
PN correction is not the right coherence test.

\section{What Velcani Argues}

Velcani used the frequency-domain expansion:
\begin{equation}
    \Psi(f)
    =
    2\pi f t_0 + \psi_0
    + \frac{3}{128\eta v^5}
      \left[1 + \sum_n v^n \psi_n \right],
    \qquad
    v=\left(\frac{\pi G M f}{c^3}\right)^{1/3},
    \qquad
    \eta=\frac{m_1m_2}{M^2}.
    \label{eq:velcani_pn_phase}
\end{equation}
Here $M=m_1+m_2$, and the $\psi_n$ are the dimensionless PN phase
coefficients. Velcani computed the 1PN correction for
$M_c = 10^{-3}M_\odot$, $f=120$ Hz and finds $v^2 \psi_2\sim \mathcal{O}(0.01)$. Higher order terms contribute 
even less because $v \sim 10^{-2} < 1$.
That is a valid expansion-order observation, but not a coherence test.

NB: Velcani uses an EOB-calibrated phenomenological model that also
allows aligned black-hole spin. However, at 1PN in the nonspinning PN limit,
the coefficient agrees with the standard nonspinning TaylorF2 PN expansion. 

\section{Computing Dephasing}
The connection between the frequency-domain PN phase and the time-domain phase $\phi(t)$
is made through the stationary phase approximation. If
$\tilde h(f)\sim A(f)e^{i\Psi(f)}$, then the stationary time associated with
frequency $f$ is
\begin{equation}
    t(f)
    =
    \frac{1}{2\pi}\frac{d\Psi}{df},
    \qquad
    \frac{dt}{df}
    =
    \frac{1}{2\pi}\frac{d^2\Psi}{df^2} .
    \label{eq:stationary_time}
\end{equation}
With this sign convention, the time-domain phase
$\phi(t)=2\pi\int^t f(t')\,dt'$ is the Legendre transform of $\Psi(f)$:
\begin{equation}
    \phi(t(f))
    =
    2\pi f\,t(f)-\Psi(f)+\mathrm{const},
    \qquad
    \phi(t)
    =
    2\pi f(t)t-\Psi(f(t))+\mathrm{const}.
    \label{eq:legendre_phase}
\end{equation}
Equation~\eqref{eq:legendre_phase} gives the phase along the chirp, but the
more direct way to accumulate phase error is to use the instantaneous
frequency. Since $d\phi/dt=2\pi f(t)$, a chirp track can be reparametrised by
frequency:
\begin{gather}
    d\phi = 2\pi f\,dt
    =
    2\pi f\,\frac{dt}{df}\,df .
\end{gather}
This is where the $dt/df$ terms and the later frequency integral come from.
The stationary-phase relation in Eq.~\eqref{eq:stationary_time} tells us how
to obtain $t(f)$ from the Fourier-domain phase $\Psi(f)$, and differentiating
it gives $dt/df$. Equivalently, once the chirp rate $df/dt$ is known, one can
use $dt/df=(df/dt)^{-1}$.

Thus PN terms in Velcani's $\Psi(f)$ expansion imply PN terms in $t(f)$, and
therefore in both $dt/df$ and the chirp rate $df/dt$. The relevant calculation
is not the fractional size of a PN coefficient at one frequency. For a coherent
resampling template, the useful quantity is the accumulated phase error after
starting the 0PN and PN tracks from the same initial frequency $f_0$. Comparing
two tracks between the same frequency endpoints gives
\begin{gather}
    \Delta\phi(f;f_0)
    =
    2\pi\int_{f_0}^{f}
    f'\left[
      \frac{dt_A}{df'} -
      \frac{dt_B}{df'}
    \right]df' ,
\end{gather}
with $A$ and $B$ later taken to be the 1PN and 0PN tracks.
Using Velcani's symbol for the nonspinning TaylorF2 coefficient, the
Fourier-domain 1PN phase term is $\psi_2v^2$, with
\begin{equation}
    \psi_2 = \frac{3715}{756}+\frac{55}{9}\eta .
    \label{eq:psi_2}
\end{equation}
Some TaylorF2 references denote this same coefficient by $\alpha_2$ rather than
$\psi_2$; this is only a change of notation for the Fourier-domain phase
$\Psi(f)$. Velcani's check evaluates the local fractional factor $\psi_2v^2$.
For the time-domain chirp, the corresponding 1PN coefficient is
\begin{equation}
    \frac{9}{20}\psi_2
    =
    \frac{743}{336}+\frac{11}{4}\eta .
    \label{eq:psi_2_time_domain}
\end{equation}
With that convention made explicit, the 1PN frequency evolution is
\begin{equation}
    \frac{df}{dt}
    =
    K f^{11/3}
    \left[
      1-\frac{9}{20}\psi_2 v^2(f)
    \right],
    \qquad
    K=\frac{96}{5}\pi^{8/3}
    \left(\frac{GM_c}{c^3}\right)^{5/3}.
    \label{eq:dfdt_1pn}
\end{equation}
Invert this to $dt/df$ and expand to first order in the 1PN term:
\begin{equation}
    \frac{dt_{\mathrm{1PN}}}{df}
    -
    \frac{dt_{\mathrm{0PN}}}{df}
    \simeq
    \frac{(9/20)\psi_2 v^2(f)}{K f^{11/3}}
    =
    \frac{9\psi_2}{20K}
    \left(\frac{\pi G M}{c^3}\right)^{2/3}
    f^{-3}.
    \label{eq:dt_df_difference}
\end{equation}
The accumulated 1PN dephasing between $f_0$ and $f$ is therefore
\begin{equation}
    \Delta\phi_{\mathrm{1PN}}(f;f_0)
    =
    2\pi\int_{f_0}^{f}
    f'\left[
      \frac{dt_{\mathrm{1PN}}}{df'} -
      \frac{dt_{\mathrm{0PN}}}{df'}
    \right]df'
    \simeq
    2\pi
    \frac{9\psi_2}{20K}
    \left(\frac{\pi G M}{c^3}\right)^{2/3}
    \left(f_0^{-1}-f^{-1}\right).
    \label{eq:dephasing_1pn}
\end{equation}
Equivalently, the cycle dephasing is
$\Delta N_{\mathrm{1PN}}=\Delta\phi_{\mathrm{1PN}}/2\pi$. For an equal-mass
$M_c=10^{-3}M_\odot$ binary started at $f_0=20\,\mathrm{Hz}$, the 1PN
accumulated dephasing by $f=120\,\mathrm{Hz}$ is
\begin{equation}
    \Delta N_{\mathrm{1PN}}(120\,\mathrm{Hz};20\,\mathrm{Hz})
    =
    2.253107577\times10^5\ \mathrm{cycles},
    \qquad
    \Delta\phi_{\mathrm{1PN}}
    =
    1.415669242\times10^6\ \mathrm{rad}.
    \label{eq:numeric_dephasing}
\end{equation}
This is the coherence-relevant number: even though
$\psi_2v^2=1.694409034\times10^{-3}$ at 120 Hz, the accumulated phase error
after starting at 20 Hz is enormous.

\end{document}
